% Chapter Template

\chapter{Ensayos y resultados} % Main chapter title

\label{Chapter4} % Change X to a consecutive number; for referencing this chapter elsewhere, use \ref{ChapterX}
Todos los capítulos deben comenzar con un breve párrafo introductorio que indique cuál es el contenido que se encontrará al leerlo.  La redacción sobre el contenido de la memoria debe hacerse en presente y todo lo referido al proyecto en pasado, siempre de modo impersonal.

%----------------------------------------------------------------------------------------
%	SECTION 1
%----------------------------------------------------------------------------------------

\section{Banco de pruebas}
Entorno utilizado para la validación: número de nodos, configuraciones, herramientas de prueba, condiciones de simulación.

\section{Pruebas unitarias}
Ejemplos de pruebas unitarias realizadas sobre los módulos principales del firmware

\section{Pruebas funcionales}
\label{sec:pruebasHW}

Pruebas de comunicación entre nodos
Pruebas de latencia y tiempo de respuesta
Pruebas de desconexión y reconexión de nodos
Pruebas de interfaz web
